\input preamble.tex
\centerline\textbf{Prøve nn  i xx }  \bigskip

Kompetansemål:
\begin{itemize}[noitemsep]
%\item utføre risikovurdering og vurdere tiltak som skal ivareta person- og maskinsikkerheten, i henhold til hvilke regelverk og normer som gjelder for arbeidet som skal utføres
%\item utføre arbeid på automatiserte anlegg fagmessig, nøyaktig og i overensstemmelse med gjeldende regelverk og normer for elektriske installasjoner og maskiner
%\item planlegge, gjennomføre, vurdere kvalitet, sluttkontrollere og digitalt dokumentere arbeidsoppdragene i faget automatiseringssystemer, individuelt og i samarbeid med andre, og begrunne valgene som er gjort
%\item tegne tekniske flytskjemaer og annen dokumentasjon, og anvende dette i utførelsen av alle arbeidsoppdrag
%\item utføre vedlikehold, systematisk feilsøking og feilretting med egnede instrumenter og verktøy, og vurdere måleresultater opp mot forventede og beregnede verdier
%\item montere, konfigurere, kalibrere, justere og sette i drift reguleringssløyfer for motor- og servodrift, temperatur, trykk, nivå og strømning, og simulere og optimalisere regulatorer basert på prosessbehov og gjøre rede for funksjon og virkemåte
%\item gjøre rede for måleprinsipper for veiing, deteksjon med kamera, vibrasjon, gassdekteksjon, pH-måling og konduktivitet for analyse og montere, konfigurere, kalibrere, justere og sette i drift minst ett av disse målesystemene
%\item montere og sette i drift ulike typer pådragsorganer med tilhørende forstillingselementer og hjelpeutstyr i henhold til leverandørens dokumentasjon, og gjøre rede for hvordan utstyret fungerer, og hvilke funksjoner det har
%\item montere og sette i drift sikkerhetskomponenter og -utstyr for nødstopp- og sikkerhetskretser, og gjøre rede for sikkerhetskategoriene Safety Integrity Level (SIL) og Performance Level (PL)
%\item programmere, montere og sette i drift programmerbare styresystemer for elektriske-, pneumatiske- og hydrauliske anlegg og gjøre rede for hvordan utstyret fungerer, og hvilke funksjoner det har
%\item installere bus-systemer og konfigurere digital og analog kommunikasjon i automatiserte anlegg, og gjøre rede for hvordan komponentene fungerer, hvordan datasignaler bearbeides og transporteres mellom nettverk, og vurdere tiltak for å etablere et sikkert elektronisk kommunikasjonsnett
%\item installere, terminere og skjøte elektriske og optiske kabler og måle kvaliteten på forbindelsen og beskrive faktorer som påvirker kvaliteten på forbindelsen
%\item installere, programmere, konfigurere og brukertilpasse «Human–Machine Interface» (HMI) for et automatiseringssystem og gjøre rede for bruksområdene
%\item simulere, sette i drift, programmere og optimalisere robot og gjøre rede for roboters funksjon og anvendelse i automatiserte anlegg
%\item foreslå, vurdere og utprøve forenklinger og forbedringer av produkter og produksjonsprosesser og gjøre rede for mulig bruk i nye produkter og tjenester
%\item håndtere avfall etter eget arbeid i industrielle og automatiserte anlegg på en miljøvennlig og økonomisk måte, drøfte produkters miljøprestasjon, slette sensitiv informasjon ved avhending og gjøre etiske og økonomiske refleksjoner rundt bærekraft, sirkulær økonomi og produktkvalitet
%\item diskutere verdien av å oppleve mestring og stolthet over eget arbeid og av å oppleve tilhørighet og trygghet i et arbeidsmiljø uavhengig av kjønn og kultur
%\item reflektere over bedriftsdemokratiets og det organiserte arbeidslivets forutsetninger, verdier og regler og hvordan et regulert arbeidsliv kan bidra til å motvirke arbeidslivskriminalitet, diskriminering og forskjellbehandling
%\item dokumentere eget arbeid, vurdere arbeidsmetoder, faglige løsninger, kvalitet og estetikk i arbeidsoppdraget, foreslå forbedringer og reflektere rundt mulige endringer
\end{itemize}

Alle ark som leveres inn skal ha elevens navn. \\ 

Oppgave 1-8 skal leveres på papir, etter levering kan eleven ta frem PC og svare på oppgave 9.\\

\vskip 2.5pt 

\hrule
\vfil \eject
Oppgave 1 (6p)%Navngi
\vskip 2.5pt 
a) \\
\vskip 2.5pt 
\begin{tikzpicture}
	\draw[step=0.5cm,gray!20,very thin]  grid (17,4) ;
\end{tikzpicture}
\vskip 2.5pt 


Oppgave 2 (6p) %Definisjoner
\vskip 2.5pt 
a) \\
\vskip 2.5pt 
\begin{tikzpicture}
	\draw[step=0.5cm,gray!20,very thin]  grid (17,4) ;
\end{tikzpicture}
\vskip 2.5pt 
Oppgave 3 (6p)%Enkel utregninger.
\vskip 2.5pt 
a) \\
\vskip 2.5pt 
\begin{tikzpicture}
	\draw[step=0.5cm,gray!20,very thin]  grid (17,4) ;
\end{tikzpicture}
\vskip 2.5pt 
Oppgave 4 (6p)
\vskip 2.5pt 
a) \\
\vskip 2.5pt 
\begin{tikzpicture}
	\draw[step=0.5cm,gray!20,very thin]  grid (17,4) ;
\end{tikzpicture}
\vskip 2.5pt 
Oppgave 5 (6p) %Feilsøkings oppgave

\vskip 2.5pt 
a) \\
\vskip 2.5pt 
\begin{tikzpicture}
	\draw[step=0.5cm,gray!20,very thin]  grid (17,4) ;
\end{tikzpicture}
\vskip 2.5pt 
\vfil \eject
Oppgave 6 (6p) % Tegn og forklar virkemåte
\vskip 2.5pt 
Tegn og forklar virkemåten til
\vfil \eject

Oppgave 7 (6p) % Tegn og forkalr virkemåte. 
\vskip 2.5pt 
Tegn og forklar virkemåten til
\vfil \eject
Oppgave 8 (6p)
\vskip 2.5pt 
a) \\
\vskip 2.5pt 
\begin{tikzpicture}
	\draw[step=0.5cm,gray!20,very thin]  grid (17,4) ;
\end{tikzpicture}
\vskip 2.5pt 
\vskip 2.5pt 
%Oppgave 9 skal beskrive hvordan en jobb skal utføres. (Planlegg, besriv hvordan du ville gjenomført og dokumenter jobben)
Oppgave 9 (12p)
\end {document}
